\newpage
\section*{第24章：存在类型}
\addcontentsline{toc}{section}{第24章：存在类型}

\begin{taplsolutionentrybox}
\phantomsection\label{sol:translator:24.3.1}
\textbf{24.3.1 解答}

先只看目标类型
\[
\forall Y.(\forall X.T\to Y)\to Y.
\]
要构造它，必须先抽象 \(Y\)，再接收
\(f:\forall X.T\to Y\)。手中只有隐藏类型 \(S\) 和值
\(t:[X\mapsto S]T\)，所以先把 \(f\) 在 \(S\) 处实例化，再应用于 \(t\)：
\[
\lambda Y.\lambda f:\forall X.T\to Y.\,f[S]t.
\]
开包时反过来，把包应用于所需结果类型 \(T_2\)，再把原来的包体写成一个能对
任意隐藏类型 \(X\) 和相应值 \(x:T_{11}\) 产生 \(T_2\) 的函数：
\[
t_1[T_2](\lambda X.\lambda x:T_{11}.t_2).
\]
这正是正文中的两条翻译。
\taplsolutionbacklink{ex:24.3.1}{返回练习24.3.1}
\end{taplsolutionentrybox}
